\section{System design}
This section discusses the application design process. The primary focus is on the requirements that must be met. A market analysis was conducted with regard to the required functions.
\subsection{Market analysis}

The market for vineyard management apps is quite niche, and most apps offer unique features. Below is a list of a few selected systems currently available on the market:

\begin{itemize}
	%\item VineInfo is a vineyard management app designed to help growers cultivate the highest-quality grapes in order to produce the best possible wine
	\item Vintrace is a cloud-based application designed to monitor and manage the entire wine production process—from the vineyard, through production and bottling, to the management of the entire supply chain. This solution ensures accessibility from anywhere, offering both a web-based application and a mobile app
	\cite{vintrace}.
	\item Vintegrate is an application designed to manage both the entire wine production process and related activities, such as direct sales to private customers or running a sommelier club
	\cite{vintegrate}.
	\item Like the solutions mentioned above, eVineyard provides systems that assist in managing the entire wine production process and also allows for the management of agricultural equipment
	\cite{evineyard}.
	\item wine.ms is another comprehensive solution that, in addition to its core features, also offers integration with select accounting systems
	\cite{wine.ms}.
	\item Mistrzwina is a Polish vineyard management system. In addition to managing the production process itself, it also assists with the registration of tax stamps and the distribution process
	\cite{mistrzwina}.
	\item Enospace is another Polish system. In addition to the features found in other competing systems, it includes a module that sends reminders about KOWR deadlines and a system for generating QR codes for bottle labels, allowing users to access information about the wines they’ve selected
	\cite{enospace}.
\end{itemize}

Current market analysis reveals that available applications are primarily designed to assist producers. However, supporting producers can be effectively integrated with increasing transparency for consumers. This solution involves creating QR codes that can be used both on equipment at the vineyard—which would help producers manage the production chain—and on bottles, which would allow consumers to trace the entire process by which their wine was made. Unlike EnoSpace’s solution, the QR code will link to a page containing not only basic information about the selected wine, but also more detailed information, such as the harvest date and location, the date the wine was placed in the barrel, and the bottling date.

\subsection{Functional requirements}
The system shall provide the following functionalities, categorized by users. This approach is consistent with the standard definition of functional requirements \cite{sommerville2017}.
\subsubsection{Guest User:}
\begin{itemize}
	\item the ability to review current offerings,
	\item the ability to recover password,
	\item the ability to log in,
	\item the ability to view information about batch of wine using qr code from the bottle including:
	\begin{itemize}
		\item final wine parameters(sugar level, acidity, pH and alcohol content),
		\item information about type of grape used for production,
		\item when was grape harvested,
		\item how long it spend in barrel if maturated,
		\item when it was bottled.
		\todo[inline]{ \textuparrow\, Tu kropka czy przecinek na końcu?}
	\end{itemize}
\end{itemize}
\subsubsection{Logged-in user:}
\begin{itemize}
	\item all of the guest user abilities,
	\item the ability to log out,
	\item the ability to change password,
	\item the ability to change email,
	\item the ability to change phone number,
	\item the ability to delete user account,
	\item the ability to create an organization and become its owner (applicable only for users with no organization affiliation),
\end{itemize}
\subsubsection{Vineyard Worker:}
\begin{itemize}
	\item all of the logged-in user abilities,
	\item the ability to record completed gardening tasks,
	\item the ability to view, edit and add plants information in specific vineyard,
	\item the ability to add data about observed pests or plant diseases,
	\item the ability to enter amount of grapes harvested from specific vineyard with particular type.
\end{itemize}
\subsubsection{Oenologist:}
\begin{itemize}
	\item all of the logged-in abilities,
	\item the ability to record chemical parameters (sugar level, acidity, pH, temperature and alcohol content) at every stage of production,
	\item the ability to enter, view, and edit data for each stage of production and bottling.
\end{itemize}
\subsubsection{Vineyard Owner:}
\begin{itemize}
	\item all of the vineyard worker and oenologist abilities,
	\item the ability to generate and review reports such as:
	\begin{itemize}
		\item production efficiency report,
		\item diversity of grapevine varieties,
		\item historical trend report.
		\todo[inline]{ \textuparrow\, Tu kropka czy przecinek na końcu?}
	\end{itemize}
	\item the ability to view and edit information about the organization,
	\item the ability to view, add and delete vineyard fields from the organization along with the grape variety planted there,
	\item the ability to view, add and delete fermenters from organization,
	\item the ability to view, add and delete users from organization,
	\item the ability to assign and revoke roles to other users in the organization.
\end{itemize}
\subsubsection{Administrator: }
\begin{itemize}
	\item all of the vineyard owner abilities,
	\item the ability to reset a user's password if they lose access to their email,
	\item the ability to change a user's first and last name.
\end{itemize}

\subsection{Non-functional requirements}

Non-functional requirements define the characteristics and requirements of a system. They are used to specify the performance parameters, technologies, and security measures on which the system is to be based. This classification provides a comprehensive understanding of the system’s operational context, going beyond its functional capabilities alone \cite{nonfunctional}.

Application described in this work should implement the following criteria:
\begin{itemize}
	\item Backend needs to be done in Java framework Spring Boot,
	\item Database have to be relational, preferably PostgreSQL,
	\item User's password should be stored in Argon2 hash,
	\item Frontend must be written in Vue.js or Svelte,
	\item Application have to work properly on Google Chrome 144 or later and Firefox 148 or later,
	\item Application needs to be accessible from mobile devices,
	\item Frontend must communicate with backend using GraphQL,
	\item All part of the system are to be conteinerized.
\end{itemize}

\subsection{User stories}
User stories are a collection of features described from the perspective of the application’s user \cite{cohn2004}. Unlike functional requirements, they provide a better understanding of the needs of non-technical users. This approach ensures that the application will be better aligned with business needs.
\subsubsection{Guest User:}
\subsubsection{Logged-in user:}

\begin{enumerate}
	\item As a logged-in user I can log out of the system so that my account is secure and no one can use it without my knowledge.
	Definition of done:
	\begin{itemize}
		\item the "log out" button is easily accessible in the user interface,
		\item after clicking the "log out" button, the user is logged out of the system, session token is invalidated on the backend,
		\item after logging out, user is redirected to the login page.
	\end{itemize}
	\item As a logged-in user I can change my password so that I can maintain the security of my account.
	Definition of done:
	\begin{itemize}
		\item the "change password" option is available in the profile settings,
		\item the user is required to provide their current password and the new password twice for confirmation,
		\item after successfully changing the password, the user receives a confirmation pop-up and is asked to log in again using the new password,
		\item the password hash is updated in the database.
	\end{itemize}
	\item As a logged-in user I can change my email address so that I can keep my personal information up to date.
	Definition of done:
	\begin{itemize}
		\item the "change email" option is available in the profile settings,
		\item the user is required to provide new email address and confirm the operation,
		\item after successfully changing the email address, the user receives a confirmation pop-up and is asked to log in again using the new email and password,
		\item the email address is updated in the database.
	\end{itemize}
	\item As a logged-in user I can change my phone number so that I can keep my personal information up to date.
	Definition of done:
	\begin{itemize}
		\item the "change phone number" option is available in the profile settings,
		\item the user is required to provide new phone number and confirm the operation,
		\item after successfully changing the phone number, the user receives a confirmation pop-up,
		\item the phone number is updated in the database.
	\end{itemize}
	\item As a logged-in user I can delete my account so that I can remove all my personal data from the system in case I no longer want to use the application.
	Definition of done:
	\begin{itemize}
		\item the "delete account" option is available in the profile settings,
		\item the user is required to enter their password and confirm the action,
		\item after successfully deleting the account, the user receives a confirmation pop-up and is redirected to the login page,
		\item all personal data is removed from the database.
	\end{itemize}
	\item As a logged-in user who is not a part of any organization I can create an organization and become its owner so that I can manage my vineyard and winery operations effectively.
	Definition of done:
	\begin{itemize}
		\item validation prevents creation if user already has an organization,
		\item the "create organization" option is available in the user dashboard,
		\item the user is required to provide the organization details and confirm the action,
		\item after successfully creating the organization, the user receives a confirmation pop-up and is redirected to the organization management page,
		\item the user is assigned the role of "owner" in the newly created organization,
		\item the organization is added to the database and associated with the user's account.
	\end{itemize}
	\item As a logged-in user I want to join an existing organization so that I can contribute to the management of the vineyard.
	Definition of done:
	\begin{itemize}
		\item the user email is visible in user dashboard so that the Vineyard Owner can easily add account to the organization,
		\item after the owner adds the user to the organization, the user receives a notification with an option to confirm,
		\item after confirming, the user is redirected to the organization dashboard and can access the features based on their assigned role,
		\item the user is associated with the organization in the database.
	\end{itemize}
\end{enumerate}

\subsubsection{Vineyard Worker:}
\subsubsection{Oenologist:}
\subsubsection{Vineyard Owner:}
\subsubsection{Administrator: }
